\section{Introduction}
\label{sec:introduction}

\subsection{The Weighted-Sum Problem}

Every redistricting optimisation tool in current use --- including every
algorithm in this platform --- employs a \emph{weighted-sum objective}:
\[
  \min\; w_1 \cdot \EC(\text{plan}) \;+\; w_2 \cdot \text{partisan\_deviation}(\text{plan})
        \;+\; w_3 \cdot \VRA(\text{plan}).
\]
The approach is convenient: it reduces a multi-criterion problem to a single
number and produces a single plan.
But it has two structural defects.

\paragraph{Defect 1: Opacity.}
The weights $w_1, w_2, w_3$ determine which trade-offs the algorithm makes,
yet their relationship to the resulting plan is non-linear and opaque.
A practitioner who wants ``compact but proportional'' plans must search the
weight space by trial and error.
The same weight vector can produce very different plans in different states
or census years; there is no principled way to calibrate weights across contexts.
In a litigation setting, opposing counsel can always argue that the weight choice
predetermined the outcome.

\paragraph{Defect 2: Hiding the trade-off structure.}
A single weighted-sum optimum hides the fact that many equally valid plans
exist with different trade-offs among the criteria.
Practitioners who want to understand ``how much compactness must we sacrifice
to achieve one more proportionally-owed seat?'' cannot answer this question
with a single-plan tool; they must run many separate optimisations with
different weight vectors and manually compare the results.

\subsection{Pareto-Optimal Redistricting}

\pareto{}-optimal redistricting addresses both defects simultaneously.
Instead of a single weighted plan, it returns a \emph{Pareto front}: a
collection of plans such that no plan in the collection is strictly better
than any other on all criteria simultaneously.
Every point on the Pareto front represents a distinct, mutually non-dominated
trade-off.
Practitioners choose from the front with full transparency: they see exactly
what compactness they sacrifice to gain proportionality, and what proportionality
they sacrifice to improve VRA compliance.

\begin{definition}[Pareto dominance]
\label{def:dominance}
Plan $A$ \emph{dominates} plan $B$ if $A$ is at least as good as $B$ on every
objective and strictly better on at least one:
\[
  A \succ B \;\iff\;
  \bigl(\forall\, m:\; A[m] \leq B[m]\bigr)
  \;\wedge\;
  \bigl(\exists\, m:\; A[m] < B[m]\bigr).
\]
The \emph{Pareto front} $\mathcal{F}_0$ is the set of all plans not dominated
by any other feasible plan.
\end{definition}

The Pareto front requires no weight choice and no linear scalarisation.
It is an objective description of the trade-off structure of the feasible space.

\subsection{Three Contributions}

This paper makes three contributions:

\begin{enumerate}
\item \textbf{NSGA-II for redistricting.}
  We adapt \nsga{}~\citep{deb2002} to the redistricting setting, using METIS
  seeds as the initial population, a ReCom-style crossover operator, and a
  single boundary-tract flip mutation.
  The result is an algorithm that maintains a diverse population of valid plans
  and converges to the Pareto front over three objectives: edge cuts (\EC),
  partisan distance ($\Dseats$), and VRA deficit ($\VRA$).

\item \textbf{Enacted-plan dominance test.}
  We introduce a binary legal test: run \nsga{} and check whether any
  Pareto-optimal plan dominates the enacted plan on all three objectives
  simultaneously.
  If so, the legislature chose a plan that is simultaneously less compact,
  less proportional, and less VRA-compliant than an algorithmically discoverable
  alternative.
  This is direct, non-statistical evidence of sub-optimal plan choice.

\item \textbf{SMC-Pareto alternative.}
  We define an alternative Pareto front construction via projection of the
  Sequential Monte Carlo (SMC) ensemble~\citep{dellaLibera2026smc} onto
  objective space.
  The SMC-Pareto projection is much faster than \nsga{} (minutes rather than
  hours) but may miss extreme regions of the Pareto front that \nsga{} actively
  searches.
\end{enumerate}

\subsection{Paper Organisation}

Section~\ref{sec:background} reviews \nsga{}~\citep{deb2002}, prior work on
multi-objective redistricting, and defines the three objectives precisely.
Section~\ref{sec:algorithm} gives the full \nsga{} algorithm with pseudocode
for all operators: non-dominated sort, crowding distance, tournament selection,
ReCom-style crossover, and boundary-tract flip mutation; it also covers the
seeding specification and the SMC-Pareto alternative.
Section~\ref{sec:results} presents empirical results on NC ($k=14$) and
WI ($k=8$): Pareto front characteristics and enacted-plan dominance tests.
Section~\ref{sec:legal} develops the legal framing of the dominance test,
compares it to ensemble-based expert arguments, and clarifies the validity
vs.\ completeness distinction in the verification protocol.
Section~\ref{sec:conclusion} positions \nsga{} redistricting in the algorithm
portfolio and identifies Phase~2 enhancements.

\paragraph{Main results.}
On NC ($k=14$, 2020 Census), \nsga{} with $N_{\mathrm{pop}}=100$ plans and
$N_{\mathrm{gen}}=200$ generations produces a Pareto front of 18--24 plans
spanning a meaningful range on all three objectives.
The enacted 2020 NC congressional map is dominated by 7 plans
on the Pareto front, meaning those plans simultaneously achieve lower edge cuts,
closer to proportional partisan balance, and better VRA compliance than the
enacted plan.
On WI ($k=8$), the enacted plan is on the Pareto frontier (not dominated),
demonstrating that the test produces negative results on some enacted plans.

\paragraph{Notation.}
$N_{\mathrm{pop}}$ denotes the NSGA-II population size (default 100);
$N_{\mathrm{gen}}$ the number of generations (default 200);
$k$ the number of congressional districts;
$n$ the number of census tracts.
A dagger~($\dagger$) on a result denotes a single run with the specified
parameters; results without a dagger are averages over multiple seeds.
